\chapter{Tiny Programs}

The first five listings are small demonstrations. Their job is not to be
complete games. Their job is to teach the rhythm of NovaBASIC and give readers
something they can finish quickly.

\section{Demo Plan}

\begin{longtable}{>{\ttfamily}p{0.18\textwidth}p{0.22\textwidth}p{0.48\textwidth}}
\toprule
File & Topic & Nova version\\
\midrule
\endhead
tiny-1.bas & Sprites & Bounce one hardware sprite with gravity, frame timing,
and sound.\\
tiny-2.bas & Sound & Play simple effects and a short musical idea with Nova's
sound commands.\\
tiny-3.bas & Line Draw & Draw points, lines, and patterns on the 320 by 200
graphics display.\\
tiny-4.bas & Vsprites & Move a larger software sprite so readers see how Nova
goes beyond 16 by 16 hardware shapes.\\
tiny-5.bas & Copper & Change display colors while the screen is being drawn.\\
\bottomrule
\end{longtable}

\section{Teaching Rule}

Each tiny program should be short enough to type in one sitting. If a feature
needs a long setup, it belongs in a later game.

\begin{typeinbox}
For printed listings, keep the first demo programs deliberately plain. A reader
should be able to type one, run it, and understand why it works without needing
the full hardware reference guide.
\end{typeinbox}

\programlesson{Tiny-1: Gravity Ball}{ball}{%
\texttt{TINY-1} puts a shiny ball on the screen, lets gravity pull it down, and
bounces it back up from the border. The program is short, but it already feels
alive because Nova's sprite hardware, frame timing, and sound commands are all
working together.}

\begin{learnbox}
\begin{itemize}
\item A hardware sprite can move without redrawing the whole screen.
\item \texttt{VSYNC} keeps the motion tied to the video frame.
\item \texttt{DY=DY+1} is enough to make simple gravity.
\item \texttt{SOUND} adds feedback without any POKEs.
\end{itemize}
\end{learnbox}

\subsection*{The Listing}

\begin{lstlisting}
50 MODE 3:COLOR 1,0,14:GCLS:GOSUB 1000
60 XL=0:XR=304:YT=0:YB=184:BV=14
65 X=40:Y=20:DX=1:DY=0
70 SPRITE 0,ON
80 SPRITE 0,X,Y:VSYNC:VSYNC
90 GET K:IF K=81 OR K=113 THEN 190
100 DY=DY+1:NX=X+DX:NY=Y+DY
120 IF NX<XL THEN NX=XL:DX=-DX:SOUND 50,2
130 IF NX>XR THEN NX=XR:DX=-DX:SOUND 50,2
140 IF NY<YT THEN NY=YT:DY=0:SOUND 60,2
150 IF NY>YB THEN NY=YB:DY=-BV:SOUND 72,3
160 X=NX:Y=NY
170 GOTO 80
190 SPRITE 0,OFF:MODE 0:COLOR 1,0,0:CLS:END
1000 REM LOAD SHINY BALL SPRITE
1010 FOR R=0 TO 15:READ A,B,C,D,E,F,H,J
1020 SPRITEDATA 0,R,A,B,C,D,E,F,H,J:NEXT R
1030 SPRITESHAPE 0,0:SPRITESET 0,6,2:SPRITESET 0,7,0:RETURN
1040 DATA $00,$00,$0F,$FF,$FF,$00,$00,$00
1050 DATA $00,$00,$FF,$FF,$FF,$FF,$00,$00
1060 DATA $00,$0F,$FF,$FF,$FF,$FF,$F0,$00
1070 DATA $00,$FF,$FF,$FF,$EE,$EE,$FF,$00
1080 DATA $0F,$FF,$FF,$EE,$EE,$EE,$EE,$F0
1090 DATA $FF,$FF,$EE,$EE,$EE,$EE,$EE,$FF
1100 DATA $FF,$FE,$EE,$EE,$EE,$EE,$EE,$6F
1110 DATA $FF,$EE,$EE,$EE,$EE,$EE,$66,$6F
1120 DATA $FE,$EE,$EE,$EE,$EE,$66,$66,$6F
1130 DATA $FE,$EE,$EE,$EE,$66,$66,$66,$BF
1140 DATA $FE,$EE,$EE,$66,$66,$66,$BB,$BF
1150 DATA $FF,$EE,$E6,$66,$66,$BB,$BB,$FF
1160 DATA $0F,$EE,$E6,$66,$BB,$BB,$BB,$F0
1170 DATA $00,$FE,$E6,$6B,$BB,$BB,$BF,$00
1180 DATA $00,$0F,$66,$BB,$BB,$BB,$F0,$00
1190 DATA $00,$0B,$BB,$BB,$BB,$BB,$B0,$00
\end{lstlisting}

\subsection*{How It Works}

Line 50 switches to graphics mode, paints the border, clears the screen, and
loads the ball sprite. Line 60 names the four walls. The screen is 320 pixels
wide and 200 pixels tall. Because the ball is 16 by 16 pixels, the right wall
is \texttt{304} and the bottom wall is \texttt{184}.

Line 100 is the whole gravity trick. Add 1 to \texttt{DY}, and the ball falls
faster. Then add \texttt{DX} to \texttt{X}, add \texttt{DY} to \texttt{Y}, and
the ball moves.

Lines 120 through 150 keep the ball inside the screen. If the next step would
cross a wall, the program puts the ball exactly on that wall. At the floor it
sets \texttt{DY=-BV}, which sends the ball upward again with the same bounce
every time. The bottom bounce uses a higher \texttt{SOUND} note than the side
bounce, so you can hear the collision without adding any POKEs.

\begin{trythisbox}
Change \texttt{BV} on line 60. A larger number makes the ball bounce higher.
A smaller number makes it bounce lower. Then change \texttt{DX} on line 65 to
make the ball drift sideways faster or slower.
\end{trythisbox}

\begin{novabox}
\texttt{SPRITE 0,X,Y} only moves the sprite registers. The picture of the ball
is already stored in sprite shape memory, so Nova does not have to redraw the
ball every frame. That is why a small BASIC loop can still look smooth.
\end{novabox}

The ball itself is also a Nova lesson. Lines 1000 through 1190 load one 16 by
16 sprite shape. The \texttt{DATA} values are the pixels. They use white, light
grey, light blue, blue, and dark grey to fake a shiny round surface.
